% sec-05: four routes that can fund a first hire.  Table 5, Figure 3
% (graphviz-type record nodes).
\section{Four Routes That Can Fund a First Hire}

An applicant's first question after "is there a role" is "is it paid for," and
the honest answer today is "not yet." This section sets out the four routes by
which a first hire could be paid, what each is worth, what gate stands in front
of it, and where it stands after today's letters. Only two of the four can pay
a salary. The third rewards jobs after they exist, and the fourth bridges a
timing gap and is never a source.

\begin{apptable}
\begin{tabularx}{\textwidth}{>{\raggedright\arraybackslash}p{3.4cm} >{\raggedright\arraybackslash}p{4.3cm} >{\raggedright\arraybackslash}X >{\raggedright\arraybackslash}p{2.6cm}}
\headrow \thead{Route} & \thead{Amount and window} & \thead{Gate in front of it} & \thead{Status today}\\
\midrule
NIH SBIR Phase I & \$306,000 over 9 months & PI employed by the company & Letter 1 sent \\
NSF SBIR Phase I & \$305,000, 6 to 18 months & A Project Pitch invitation & Pitch on day 4 \\
California Competes & A negotiated tax credit & New California jobs & Letter 2 asks \\
Company reserve & \$40,000 for 13 weeks & Award letter and offer & Rung A named \\
\end{tabularx}
\end{apptable}
\vspace{-0.60cm}
\tabcap{Table 5. Four routes that could pay for a first hire, with what each is worth, the\\
gate in front of it, and where it stands; only the first two can pay a salary outright.}

The two federal routes are close in size and different in shape. The NIH Phase
I is the company's own planned budget of \$306,000, carried unchanged from its
ten-application work \cite{tenapps2026,nihsbirseed}. The NSF Phase I ceiling
under the current solicitation is \$305,000, and NSF will not take a full
proposal without first inviting one through a Project Pitch \cite{nsf26510}. A
consenting applicant could be named in either, in the form the NIH letter asks
about, and never in both for the same effort.

\begin{appfloat}[!tb]
\begin{appfig}[graphviz-type / record nodes under one root]
% One root ellipse and four record nodes, each a column of five ruled cells:
% the route, then amount, window, gate and status.  A gray field-name column
% on the left labels the rows once, as a graphviz record's field ports would.
\node[gvkey,text width=30mm] (root) at (5.35,1.55) {A first hire, paid on award};
\foreach \y/\f in {0/Route, -0.72/Amount, -1.44/Window, -2.16/Gate, -2.88/Today}{
  \node[gvcellg,text width=17mm,minimum height=7mm] at (-2.35,\y) {\f};}
\foreach \x/\a/\b/\c/\d/\e/\h in {%
  0.00/{NIH SBIR Phase I}/{\$306,000 total}/{9 months}/{PI primarily employed by the company}/{Letter 1 sent}/gvcellh,
  3.55/{NSF SBIR Phase I}/{Up to \$305,000}/{6 to 18 months}/{Project Pitch invitation first}/{Pitch on day 4}/gvcellh,
  7.10/{California Competes}/{Negotiated tax credit}/{GO-Biz periods}/{California jobs created}/{Letter 2 asks}/gvcellh,
  10.65/{Company reserve}/{\$40,000 ceiling}/{90 days pre-award}/{Award letter and offer}/{Rung A designated}/gvcells}{
  \node[\h,text width=30mm,minimum height=7mm] at (\x,0) {\a};
  \node[gvcell,text width=30mm,minimum height=7mm] at (\x,-0.72) {\b};
  \node[gvcell,text width=30mm,minimum height=7mm] at (\x,-1.44) {\c};
  \node[gvcell,text width=30mm,minimum height=7mm] at (\x,-2.16) {\d};
  \node[gvcell,text width=30mm,minimum height=7mm] at (\x,-2.88) {\e};}
\draw[gvedgeb] (root.south west) to[out=-120,in=90] (0.00,0.36);
\draw[gvedgeb] (root.south) to[out=-100,in=90] (3.55,0.36);
\draw[gvedge]  (root.south) to[out=-80,in=90] (7.10,0.36);
\draw[gvedged] (root.south east) to[out=-60,in=90]
  node[pos=0.55,right,font=\tiny\itshape,fill=fundwhite,inner sep=1pt] {bridge only} (10.65,0.36);
\ptitle{-3.25}{1.55}{Four routes, five fields}
\pnote{-3.25}{-3.72}{Bold edges pay salary; the plain edge credits tax once jobs exist; the dashed edge only bridges timing.}
\end{appfig}
\vspace{-0.60cm}
\figcaption{Figure 3. The four routes as records under one root, with the same five fields on\\
each, so that amount, window, gate and status can be read straight across the rows.}
\end{appfloat}

\subsection{The money frame a first hire sits inside}

None of the numbers below is re-derived for this packet. The program's direct
cost is \$700,000 a year for five years, \$3,500,000 in total, of which
\$521,000 a year is personnel across 3.95 full-time equivalents
\cite{movein2026}. The SBIR route through Phase I at \$306,000 and Phase II at
\$1,300,000 totals \$1,606,000, leaving a delta of \$2,104,000 that the SBIR
route does not buy; the capitalization plan places \$5,900,000 of private
capital behind a firewall against that delta, a private to federal leverage of
3.67 to 1 \cite{capitalplan2026}. The company's comparable virtual trial work is
\textbf{projected} at \$36,330 per run, and that figure describes simulation work
already deposited rather than the cost of the clinical trial.

\subsection{What a funder buys when it pays for these roles}

A funder who pays for the six open roles is paying for the people who would run
a specific trial: open-label, single-arm, 3+3 escalation of perioperative
daraxonrasib at 160, 220 and 300 mg, up to 18 participants, a 28-day
dose-limiting toxicity window, and a staged eight-arm robotic
pancreaticoduodenectomy with 56 degrees of freedom, 640 sensor channels and a
five-vessel no-fly gate, advised but never controlled by an on-premises language
model \cite{onpremwhippl,phase1ind}. The evidence behind the agent choice is
dated and public: a 2025 simulation that returned 12.8 against 5.4 months
\cite{qspmetpancre}, and a 2026 Phase 3 readout of 13.2 against 6.6 months in
the RAS G12 population \cite{rasolute302}. That proximity is a chronology
observation, not a validation claim: the simulation modeled 1000 patients and a
combination, and the trial enrolled a different population on a single agent.

Pancreatic ductal adenocarcinoma remains the disease with the widest gap between
incidence and survival in the United States \cite{siegel2025statistics}, and the
standard of care after resection is still combination chemotherapy
\cite{conroy2018folfirinox}. The developer's own adjuvant Phase 3 in resected
disease, RASolute 304, is now enrolling \cite{ctgov304}. The perioperative,
robot-assisted use this program proposes is a different question from either
trial, and it is the question the six roles would be hired to answer.
